% TOP_PROBLEM: {"id":30007547,"problem_number":"LOCAL-30007547","title":"Entrant versus incumbent innovation","category_id":21,"category":{"id":21,"name":"economics","display_name":"Economics","description":"Open economic theory statements, published research agendas, and proposed scoped specifications; question provenance and historical priority are qualified per record.","slug":"economics","order_index":21,"created_at":"2026-10-08T00:00:00Z"},"status":"research_question","external_url":null,"legacy_ids":[],"published":false,"statement":"Identify the causal response of incumbent own-product research to entrant innovation hazard, separating partial and general-equilibrium responses.","economics":{"original_id":36,"field":"Growth and productivity","kind":"identification","formalization_basis":"adapted_research_question","source_status":"explicit_open","priority_status":"earliest_found","priority_note":"This is the earliest explicit relevant statement verified in this audit, not a claim of original priority; older literature and earlier versions may contain antecedents.","earliest_verified_reference":{"authors":["Peter J. Klenow","Huiyu Li"],"title":"Innovative Growth Accounting","year":2021,"url":"https://www.journals.uchicago.edu/doi/full/10.1086/712325","doi":"10.1086/712325","locator":"Section VII, Conclusion, penultimate paragraph on determinants of innovation arrival rates","evidence_summary":"The authors explicitly identify as a future research question whether creative destruction is a strategic complement or substitute for incumbent own-product innovation. Their accounting decomposition does not itself answer that strategic question."},"current_reference":{"authors":["Peter J. Klenow","Huiyu Li"],"title":"Innovative Growth Accounting","year":2021,"url":"https://www.journals.uchicago.edu/doi/full/10.1086/712325","doi":"10.1086/712325","locator":"Section VII, Conclusion, penultimate paragraph on determinants of innovation arrival rates","evidence_summary":"The authors explicitly identify as a future research question whether creative destruction is a strategic complement or substitute for incumbent own-product innovation. Their accounting decomposition does not itself answer that strategic question."},"formalization_note":"The conclusion explicitly names this complement/substitute question. The response derivative and identification design are adapted, with partial and general-equilibrium effects separated."},"statusQualification":"Published question or research agenda verified at the cited locator. The mathematical specification is adapted for this catalogue and includes additional assumptions; source publication does not establish first-ever priority or certify this exact model remains unsolved. The conclusion explicitly names this complement/substitute question. The response derivative and identification design are adapted, with partial and general-equilibrium effects separated.","statusReviewedAt":"2026-10-08"}
% FIRST_POSTED: 2026-10-08
% ENTRY_KIND: statement_only
% !TeX program = lualatex
\documentclass[11pt]{article}
\usepackage{fontspec}
\usepackage[a4paper,margin=25mm]{geometry}
\usepackage{amsmath,amssymb,mathtools}
\usepackage{xurl}
\usepackage[hidelinks]{hyperref}
\newcommand{\sref}[2]{\hyperlink{source:#1:#2}{[S#2]}}
\setlength{\emergencystretch}{3em}
\begin{document}
% problemId:30007547
\section*{Entrant versus incumbent innovation}\label{30007547}
\subsection{Definitions and mathematical statement}
Fix a product-line innovation model with incumbents choosing research intensity \(r\ge0\) and potential entrants generating creative-destruction hazard \(\lambda\ge0\). Let \(V(r,\lambda;w,p)\) be incumbent expected discounted profit net of research cost, with wage \(w\) and demand parameters \(p\) initially fixed. Under a unique differentiable interior best response \(r^*(\lambda;w,p)\), define strategic complementarity or substitutability by the sign of
\[
\frac{\partial r^*(\lambda;w,p)}{\partial\lambda}.
\]
Identify this derivative over a declared range of entrant pressure and distinguish it from the general-equilibrium derivative when research wages, product demand, and financing costs also change. Specify incumbent own-product innovation separately from acquisition of entrants or innovation in unrelated lines.
Use instruments shifting entrant innovation opportunities while satisfying stated exclusions for incumbent demand and research productivity; provide bounds if those exclusions are only partly defensible. Match incumbent research, product-quality changes, and entrant hazard, rather than inferring strategic effects from aggregate entry alone. Characterize mechanisms under which replacement risk discourages research or an escape-from-competition incentive raises it, and test whether their quantitative balance varies across technological distance. The target is a causal strategic-response function, not merely an accounting decomposition of growth into incumbent and entrant contributions.

\par\textbf{Formulation and scope.} The conclusion explicitly names this complement/substitute question. The response derivative and identification design are adapted, with partial and general-equilibrium effects separated.

\par\textbf{Open-question provenance.} An explicit question or research agenda was verified at the source locator. This does not by itself certify that every later version remains unresolved \sref{30007547}{1}.

\par\textbf{Historical priority.} This is the earliest explicit relevant statement verified in this audit, not a claim of original priority; older literature and earlier versions may contain antecedents.

\subsection{Short English statement}
Identify the causal response of incumbent own-product research to entrant innovation hazard, separating partial and general-equilibrium responses.

\subsection{Sources}
\begin{itemize}\raggedright
\item[\textnormal{[S1]}]\hypertarget{source:30007547:1}{} Peter J. Klenow, Huiyu Li. 2021. Innovative Growth Accounting. Section VII, Conclusion, penultimate paragraph on determinants of innovation arrival rates.
\url{https://www.journals.uchicago.edu/doi/full/10.1086/712325}.
DOI: 10.1086/712325.
{\small\textit{Earliest verified question/agenda reference; historical priority qualified below. The authors explicitly identify as a future research question whether creative destruction is a strategic complement or substitute for incumbent own-product innovation. Their accounting decomposition does not itself answer that strategic question.}}
\end{itemize}
\end{document}
